\documentclass{grattan}
% Comments are deployed by the % sign; everything after % is ignored by the compiler.
% Please do not put comments before \documentclass as these are reserved for TeX directives.

\author{}
\title{}

\GrattanReportNumber{2017-00}

\acknowledgements{%
This 
%
}


\begin{document}

\begin{overview}
This is the overview. 
The word \textbf{Overview} appears in the contentspage at \verb=\contentspage= as an unnumbered entry.
Set over two columns, the first and last baseline of each column is set on matching baselines. 

(But only if the overview is large enough!) 

The first line of the overview starts at a fixed vertical position but its height is determined by the constraint of matching the baselines.
In particular if the overview is too large, it will overflow the page.

Footnotes are not permitted.
\end{overview}

\begin{recommendations}
\recommendation{Recommendation}
Some text under a recommendation.
\end{recommendations}

\contentspage

\chapter{Chapter title}
\section{Section title}
\subsection{Subsection title}
\subsubsection{Subsubsection title}
Some text in one paragraph. 
Some text in one paragraph. 
Some text in one paragraph. 

Some text in the next paragraph. 
Some text in the next paragraph. 
Some text in the next paragraph. 

Since singular newlines are ignored, put different sentences on different lines.

\begin{enumerate}
	\item An enumerated list
	\item Another item
	\begin{enumerate}
		\item A subitem of the second item
		\item Another subitem

		New paragraphs are the same as in the body text
	\end{enumerate}
\end{enumerate}

Another list:
\begin{itemize}
	\item A bulleted/non-enumerated list
	\item Another item
	\begin{itemize}
		\item A subitem.
	\end{itemize}
\end{itemize}

\chapter{Figures and tables}
An initial cross-reference to \Vref{fig:first-example-figure}.

\begin{figure}
\caption{Figure caption\label{fig:first-example-figure}}
\units{Figure units}
\includegraphics{logos/GrattanSVGLogo.pdf}
\noteswithsource{Notes}{Source}
\end{figure}

A cross-reference to \Vref{box:example}.

\begin{verysmallbox}{A very small box}{box:very-small-box-example}
Contents of very smallbox.

A paragraph with a footnote.

Some text in a very small box. 
Some text in a very small box. 
Some text in a very small box. 
\end{verysmallbox}

\begin{smallbox}{A small box}{box:example}
A smallbox. 
Some text in a small box. 
Some text in a small box. 
Some text in a small box. 
Some text in a small box.\footnote{Footnote entry.}


Some text in a small box. 
Some text in a small box. 
Some other text in a small box. 
Some other text in a small box. 
Some text in a small box. 
Some text in a small box. 
Some text in a small box. 
Some other text in a small box. 
Some other text in a small box. 
Some text in a small box. 


Some text in a small box. 
Some text in a small box. 
Some text in a small box. 
Some text in a small box. 
Some text in a small box.

Some text in a small box. 
Some text in a small box. 
Some other text in a small box. 
Some other text in a small box. 
Some text in a small box. 

Some text in a small box. 
Some text in a small box. 
Some other text in a small box. 
Some other text in a small box. 
Some text in a small box. 
Some text in a small box. 
Some text in a small box. 
Some other text in a small box. 
Some other text in a small box. 
Some text in a small box. 

Some text.\footnote{All footnotes should end with a full stop.} 
Some other text.\footnote{(Or a closing parenthesis, which itself should be preceded by a full stop.)}
\end{smallbox}

A cross-reference to \Vref{fig:2nd-example-figure}, with a reference \Cref{fig:first-example-figure}.

\begin{figure}
\caption{A second figure caption\label{fig:2nd-example-figure}}
\units{Figure units}
\includegraphics{logos/GrattanSVGLogo.pdf}
\noteswithsource{Notes}{Source}
\end{figure}

A table:

\begin{table}
\caption{Table caption}\label{tbl:one-table}
\begin{tabularx}{\linewidth}{XXR}
%
\toprule
\textbf{Column title}                           & \textbf{Column type}                              & \textbf{Ragged left title} \\
\midrule
Move from one cell to the next cell in the same row by using an ampersand & Move to the next line by using a double backslash & Use toprule, midrule, and bottomrule. \\
Move from one cell to the next cell in the same row by using an ampersand & Move to the next line by using a double backslash & Use toprule, midrule, and bottomrule. \\
Move from one cell to the next cell in the same row by using an ampersand & Move to the next line by using a double backslash & Use toprule, midrule, and bottomrule. \\
Move from one cell to the next cell in the same row by using an ampersand & Move to the next line by using a double backslash & Use toprule, midrule, and bottomrule. \\[15.5pt]
Use the optional argument & to the double backslash & to specify the precise distance between rows. \\
\cmidrule(lr){2-3}
Never use vertical rules & But sometimes horizontal rules are appropriate. & Use cmidrule(lr) \\
\bottomrule
\end{tabularx}
\noteswithsource{Notes}{Source}



\end{table}
Another table, cross-referenced in the same way \Vref{tbl:one-table}.


\end{document}
